\chapter{Tek Yönlü Varyans Analizi: Üç veya Daha Fazla Ortalamanın Karşılaştırılması}
\label{ch:anova}

\begin{prismbox}[PrismLight]{Öğrenme çıktıları}
Bu bölümün sonunda:
\begin{itemize}
  \item tek yönlü bağımsız gruplar ANOVA'sının araştırma sorusunu ve hipotezlerini kurabilecek,
  \item toplam değişkenliği gruplar arası ve grup içi bileşenlere ayırabilecek,
  \item ANOVA tablosundaki kareler toplamı, serbestlik derecesi, ortalama kare ve $F$ değerini yorumlayabilecek,
  \item genel testten sonra uygun çoklu karşılaştırmaları ve etki büyüklüklerini değerlendirebilecek,
  \item klasik ANOVA ile Welch ANOVA arasında veri yapısına göre seçim yapabileceksiniz.
\end{itemize}
\end{prismbox}

\begin{prerequisitebox}
Bağımsız örneklem \tstat testi, hipotez testi, güven aralığı, etki büyüklüğü
ve grup içi değişkenlik kavramları bilinmelidir. İki bağımsız grubun
karşılaştırılması için Bölüm~\ref{ch:two-sample}, genel test mantığı için
Bölüm~\ref{ch:hypothesis} temel alınır.
\end{prerequisitebox}

\section{Neden ANOVA?}

Üç veya daha fazla bağımsız grubun ortalamaları karşılaştırılırken bütün
ikili \tstat testlerini ayrı ayrı yapmak aile düzeyindeki Tip I hata riskini
artırır. Tek yönlü varyans analizi (ANOVA), grupların ortalamalarının aynı
olduğu ortak bir sıfır hipotezini tek bir genel testle sınar. Yöntemin
hesaplama ve raporlama ilkeleri için bkz. \citet{agresti2018}; etki
büyüklüklerinin yorumu için ayrıca \citet{lakens2013}.

Tek yönlü ANOVA'da bir kategorik açıklayıcı değişken ve bir nicel yanıt
değişkeni vardır. Örneğin üç öğretim yöntemindeki başarı puanları, dört okul
türündeki aidiyet puanları veya üç geri bildirim biçimindeki tamamlama süreleri
karşılaştırılabilir. Aynı bireylerin üç zamanda ölçülmesi ise bağımsız gruplar
ANOVA'sı değildir; tekrarlı ölçümler için bağımlılığı dikkate alan yöntem gerekir.

\begin{definitionbox}
\textbf{Tek yönlü bağımsız gruplar ANOVA'sı}, bir nicel yanıtın $k$ bağımsız
gruptaki evren ortalamalarını karşılaştırır. ``Tek yönlü'' ifadesi tek bir
gruplama değişkeni bulunduğunu; testin tek kuyruklu olduğunu değil, $F$
dağılımının sağ kuyruğunun kullanıldığını belirtir.
\end{definitionbox}

\section{Hipotezler ve genel testin anlamı}
\label{sec:anova-hypotheses}

$k$ grubun evren ortalamaları $\mu_1,\ldots,\mu_k$ olmak üzere hipotezler
\[
H_0:\mu_1=\mu_2=\cdots=\mu_k,
\qquad
H_1:\text{en az bir evren ortalaması farklıdır}
\]
biçimindedir. Alternatif hipotez bütün ortalamaların birbirinden farklı
olduğunu söylemez. Anlamlı genel test yalnız eşitlik modelinin veriyle uyumlu
olmadığına işaret eder; farkın hangi gruplar arasında olduğunu belirlemek için
planlı karşılaştırma veya çoklu karşılaştırma gerekir.

\begin{mistakebox}
Anlamlı $F$ testi ``bütün gruplar farklıdır'' anlamına gelmez. Benzer biçimde
anlamlı olmayan genel test, bütün ortalamaların kesin olarak eşit olduğunu
kanıtlamaz; güven aralıkları ve çalışmanın hassasiyeti ayrıca incelenmelidir.
\end{mistakebox}

\section{Değişkenliğin iki kaynağa ayrılması}
\label{sec:anova-decomposition}

ANOVA, puanların genel ortalama çevresindeki toplam değişkenliğini iki parçaya
ayırır:
\[
SS_{\text{toplam}}=SS_{\text{gruplar}}+SS_{\text{hata}}.
\]
Burada $SS$ kareler toplamını gösterir. Grup $j$'nin hacmi $n_j$, ortalaması
$\bar x_j$ ve genel ortalama $\bar x_{..}$ ise
\[
SS_{\text{gruplar}}
=\sum_{j=1}^{k}n_j(\bar x_j-\bar x_{..})^2
\]
grup ortalamalarının genel ortalamadan uzaklığını ölçer. Grup içindeki
gözlemler $x_{ij}$ ile gösterildiğinde
\[
SS_{\text{hata}}
=\sum_{j=1}^{k}\sum_{i=1}^{n_j}(x_{ij}-\bar x_j)^2
\]
aynı gruptaki bireysel farklılıkları özetler.

\begin{figure}[H]
\centering
\begin{tikzpicture}[x=1.05cm,y=0.16cm]
  \draw[->,PrismGray] (0.4,0)--(9.3,0) node[right,font=\scriptsize]{puan};
  \foreach \x/\lab in {1/60,3/65,5/70,7/75,9/80}
    \node[below,font=\tiny] at (\x,0) {\lab};
  \foreach \x/\y in {1.8/3,2.4/5,3.0/3.5,3.5/5.4,4.0/3.8}
    \fill[PrismBlue] (\x,\y) circle (2pt);
  \foreach \x/\y in {3.7/9,4.5/10.8,5.1/8.7,5.8/10.1,6.3/9.2}
    \fill[PrismSubheadingBlue] (\x,\y) circle (2pt);
  \foreach \x/\y in {5.8/14.3,6.5/15.8,7.2/13.8,7.9/15.3,8.5/14.5}
    \fill[PrismWarnOrange!85!black] (\x,\y) circle (2pt);
  \draw[PrismBlue,very thick] (3.0,2)--(3.0,6.2);
  \draw[PrismSubheadingBlue,very thick] (5.1,7.7)--(5.1,11.8);
  \draw[PrismWarnOrange!85!black,very thick] (7.2,12.8)--(7.2,16.7);
  \node[left,font=\scriptsize] at (0.7,4.5) {A};
  \node[left,font=\scriptsize] at (0.7,9.8) {B};
  \node[left,font=\scriptsize] at (0.7,14.8) {C};
  \draw[PrismGray,dashed] (5.1,0)--(5.1,17.2)
    node[above,font=\scriptsize]{genel ortalama};
\end{tikzpicture}
\caption{ANOVA'da grup ortalamaları arasındaki uzaklık ile grup içi yayılımın birlikte değerlendirilmesi.}
\label{fig:anova-variation}
\end{figure}

Şekil~\ref{fig:anova-variation}'de dikey kısa çizgiler grup ortalamalarını,
noktaların bu çizgiler çevresindeki saçılması grup içi değişkenliği gösterir.
Grup ortalamaları birbirinden uzak ve grup içi yayılım küçük olduğunda genel
$F$ oranı büyür.

\section[ANOVA tablosu ve F istatistiği]{ANOVA tablosu ve $F$ istatistiği}
\label{sec:anova-table}

Kareler toplamları serbestlik derecelerine bölünerek ortalama kareler elde edilir:
\[
MS_{\text{gruplar}}=\frac{SS_{\text{gruplar}}}{k-1},
\qquad
MS_{\text{hata}}=\frac{SS_{\text{hata}}}{N-k}.
\]
Genel test istatistiği
\[
F=\frac{MS_{\text{gruplar}}}{MS_{\text{hata}}}
\]
olup $H_0$ altında yaklaşık $F_{k-1,N-k}$ dağılımını izler. Eşit ortalamalar
altında pay ve payda aynı hata değişkenliğini tahmin ettiğinden oran yaklaşık
1 çevresindedir. Grup ortalamaları arasındaki fark büyüdükçe pay ve dolayısıyla
$F$ değeri artar. Kritik değerler Ek~\ref{app:distribution-tables},
Bölüm~\ref{sec:f-table}'de verilmiştir.

\begin{table}[H]
\centering
\caption{Tek yönlü ANOVA tablosunun genel yapısı.}
\label{tab:anova-template}
\begin{tabularx}{\textwidth}{@{}>{\raggedright\arraybackslash}Xrrrr@{}}
\toprule
Kaynak & Kareler toplamı & sd & Ortalama kare & $F$ \\
\midrule
Gruplar arası & $SS_{\text{gruplar}}$ & $k-1$ & $MS_{\text{gruplar}}$ & $MS_{\text{gruplar}}/MS_{\text{hata}}$ \\
Grup içi (hata) & $SS_{\text{hata}}$ & $N-k$ & $MS_{\text{hata}}$ & \\
Toplam & $SS_{\text{toplam}}$ & $N-1$ & & \\
\bottomrule
\end{tabularx}
\end{table}

\section{Çözümlü uygulama: üç öğretim yöntemi}
\label{sec:anova-worked-example}

Üç bağımsız öğretim yönteminde onar öğrencinin başarı puanı özetleri şöyle
olsun:

\begin{table}[H]
\centering
\caption{Üç öğretim yönteminin başarı puanı özetleri.}
\label{tab:anova-group-summary}
\begin{tabular}{lrrr}
\toprule
Yöntem & $n$ & Ortalama & Standart sapma \\
\midrule
A & 10 & 70 & 6 \\
B & 10 & 75 & 7 \\
C & 10 & 82 & 5 \\
\bottomrule
\end{tabular}
\end{table}

Genel ortalama
\[
\bar x_{..}=\frac{10(70)+10(75)+10(82)}{30}=75{,}667
\]
olur. Gruplar arası kareler toplamı
\[
SS_{\text{gruplar}}
=10(70-75{,}667)^2+10(75-75{,}667)^2+10(82-75{,}667)^2
=726{,}667
\]
ve özet standart sapmalardan grup içi kareler toplamı
\[
SS_{\text{hata}}
=(10-1)6^2+(10-1)7^2+(10-1)5^2=990
\]
bulunur. Sonuç Tablo~\ref{tab:anova-worked}'de gösterilmiştir.

\begin{table}[H]
\centering
\caption{Üç öğretim yöntemi için ANOVA tablosu.}
\label{tab:anova-worked}
\begin{tabular}{lrrrrr}
\toprule
Kaynak & $SS$ & sd & $MS$ & $F$ & $p$ \\
\midrule
Gruplar arası & 726,667 & 2 & 363,333 & 9,909 & $<0{,}001$ \\
Grup içi & 990,000 & 27 & 36,667 & & \\
Toplam & 1716,667 & 29 & & & \\
\bottomrule
\end{tabular}
\end{table}

$F(2,27)=9{,}91$ ve $p=0{,}00059$ olduğundan eşit ortalamalar hipotezi
reddedilir. Sonuç en az iki yöntemin evren ortalamalarının farklı olduğuna
ilişkin kanıt sağlar; hangi çiftlerin ayrıldığını henüz göstermez.

\section{Etki büyüklüğü}
\label{sec:anova-effect-size}

Genel ilişkinin örneklemdeki büyüklüğü eta-kare ile
\[
\eta^2=\frac{SS_{\text{gruplar}}}{SS_{\text{toplam}}}
\]
özetlenebilir. Eta-kare örneklemde yanıt değişkenliğinin gruplama değişkeniyle
ilişkili bölümüdür. Özellikle küçük örneklemde yukarı yönlü yanlılığı azaltmaya
çalışan omega-kare
\[
\omega^2=
\frac{SS_{\text{gruplar}}-(k-1)MS_{\text{hata}}}
{SS_{\text{toplam}}+MS_{\text{hata}}}
\]
ile hesaplanabilir.

Çözümlü örnekte
\[
\eta^2=\frac{726{,}667}{1716{,}667}=0{,}423,
\qquad
\omega^2\approx0{,}373.
\]
Örneklemde başarı puanı değişkenliğinin yaklaşık yüzde 42,3'ü öğretim yöntemi
gruplarıyla ilişkilidir. Bu oran nedensel etki değildir; nedensel yorum için
öğrencilerin yöntemlere atanma süreci belirleyicidir. Küçük, orta ve büyük
gibi sabit etiketler yerine alan bilgisi, ölçme birimi ve güven aralığıyla
birlikte yorum yapılmalıdır.

\section{Genel testten sonra çoklu karşılaştırmalar}
\label{sec:anova-posthoc}

Genel test anlamlı olduğunda araştırma sorusu hangi grup farklarının önemli
olduğunu gerektiriyorsa aile düzeyindeki hata oranını denetleyen bir yöntem
kullanılır. Eşit varyanslı bağımsız gruplarda Tukey HSD bütün ikili farklar
için yaygın bir seçenektir. Yalnız önceden belirlenmiş birkaç bilimsel
karşılaştırma varsa planlı karşıtlıklar daha doğrudan olabilir.

Çözümlü örneğin Tukey sonuçları şöyledir:

\begin{table}[H]
\centering
\caption{Öğretim yöntemleri için Tukey HSD karşılaştırmaları.}
\label{tab:anova-tukey}
\begin{tabular}{lrrr}
\toprule
Karşılaştırma & Ortalama farkı & Yüzde 95 eşzamanlı GA & Düzeltilmiş $p$ \\
\midrule
A $-$ B & $-5$ & $[-11{,}71;1{,}71]$ & 0,174 \\
A $-$ C & $-12$ & $[-18{,}71;-5{,}29]$ & $<0{,}001$ \\
B $-$ C & $-7$ & $[-13{,}71;-0{,}29]$ & 0,040 \\
\bottomrule
\end{tabular}
\end{table}

Tablo~\ref{tab:anova-tukey}'ye göre C yöntemi A ve B yöntemlerinden daha
yüksek ortalamaya sahiptir; A ile B arasındaki beş puanlık fark için yeterli
kanıt yoktur. Düzeltilmiş $p$ değerleri ve eşzamanlı aralıklar aynı karşılaştırma
ailesini dikkate alır.

\begin{mistakebox}
Önce genel ANOVA yapıp ardından düzeltmesiz çok sayıda \tstat testi uygulamak
çoklu karşılaştırma sorununu çözmez. Karşılaştırmaların veri görüldükten sonra
seçilmesi de doğrulayıcı ve keşfedici analiz ayrımını bulanıklaştırır.
\end{mistakebox}

\section{Varsayımlar ve sağlam seçenekler}
\label{sec:anova-assumptions}

Klasik tek yönlü ANOVA için temel noktalar şunlardır:

\begin{itemize}
  \item Gözlemler araştırma tasarımına göre bağımsız olmalıdır.
  \item Her gruptaki hata dağılımı ağır aykırı değerlerden ve ciddi biçim bozukluklarından uzak olmalıdır.
  \item Grup hata varyanslarının yaklaşık eşit olduğu varsayılır.
  \item Yanıt nicel, gruplar birbirini dışlayan kategoriler olmalıdır.
\end{itemize}

Dengeli örneklem hacimleri ve benzer dağılımlar klasik ANOVA'yı orta düzey
normallik sapmalarına karşı dayanıklı kılar. Varyanslar ve grup hacimleri
belirgin biçimde farklıysa, özellikle küçük grubun varyansı büyükse, Welch
ANOVA daha güvenli bir genel test sunar. Welch testi eşit varyans varsayımını
gerektirmez ve payda serbestlik derecesini yaklaşık hesaplar. Böyle bir
analizden sonra Games--Howell gibi eşit varyans gerektirmeyen karşılaştırmalar
düşünülmelidir.

Sıralara dayalı Kruskal--Wallis testi dağılımlar arasındaki genel farklılığı
inceler; ek koşullar olmadan otomatik olarak ``medyanların eşitliği testi''
diye yorumlanmamalıdır. Bağımsızlığın bozulması ise dönüşüm veya sıralama
yöntemiyle düzeltilemez; tasarıma uygun model gerekir.

\begin{assumptionbox}
Levene testinin anlamsız çıkması varyansların eşit olduğunu kanıtlamaz;
anlamlı çıkması da tek başına veri dönüşümü zorunluluğu oluşturmaz. Grup
hacimleri, standart sapmalar, kutu grafikleri, aykırı değerler ve klasik--Welch
duyarlılık karşılaştırması birlikte değerlendirilmelidir.
\end{assumptionbox}

\section{Welch ANOVA örneği}

Üç grubun özetleri sırasıyla $(n,\bar x,s)=(12,70,3)$,
$(20,76,10)$ ve $(8,84,5)$ olsun. Varyanslar 9, 100 ve 25 olduğundan eşit
varyans varsayımı zayıftır; grup hacimleri de dengesizdir. Welch hesabı
yaklaşık
\[
F_W(2,18{,}16)=25{,}33,\qquad p<0{,}001
\]
sonucunu verir. Ortalamaların eşit olmadığına ilişkin güçlü kanıt vardır.
Hangi grupların ayrıldığını belirlemek için eşit varyans gerektirmeyen ikili
karşılaştırmalar ve fark güven aralıkları kullanılmalıdır.

\section{Yazılım uygulaması}

\begin{softwarebox}
\begin{lstlisting}
import numpy as np
from scipy import stats
from statsmodels.stats.oneway import anova_generic

# Ozet istatistiklerden klasik ANOVA
n = np.array([10, 10, 10])
mean = np.array([70., 75., 82.])
sd = np.array([6., 7., 5.])
grand = np.sum(n * mean) / np.sum(n)
ss_between = np.sum(n * (mean - grand)**2)
ss_within = np.sum((n - 1) * sd**2)
df_between, df_within = len(n) - 1, np.sum(n) - len(n)
mse = ss_within / df_within
f_value = (ss_between / df_between) / mse
p_value = stats.f.sf(f_value, df_between, df_within)
eta2 = ss_between / (ss_between + ss_within)
omega2 = ((ss_between - df_between * mse) /
          (ss_between + ss_within + mse))

print(f_value, p_value, eta2, omega2)

# Ozetlerden Tukey HSD eszamanli araliklari
q_critical = stats.studentized_range.ppf(
    .95, len(n), df_within)
se_tukey = np.sqrt(mse / n[0])
for i, j in [(0, 1), (0, 2), (1, 2)]:
    diff = mean[i] - mean[j]
    margin = q_critical * se_tukey
    p_adjusted = stats.studentized_range.sf(
        abs(diff) / se_tukey, len(n), df_within)
    print(i, j, diff, (diff-margin, diff+margin),
          p_adjusted)

# Esit olmayan varyanslarda Welch ANOVA: ozetlerden
welch = anova_generic(
    means=np.array([70., 76., 84.]),
    variances=np.array([3., 10., 5.])**2,
    nobs=np.array([12, 20, 8]),
    use_var="unequal", welch_correction=True)
print(welch.statistic, welch.pvalue, welch.df)
\end{lstlisting}
İlk hesap yaklaşık $F(2,27)=9{,}91$, $p=0{,}00059$,
$\eta^2=0{,}423$ ve $\omega^2=0{,}373$ verir. Döngü,
Tablo~\ref{tab:anova-tukey}'deki Tukey farklarını, eşzamanlı aralıkları ve
düzeltilmiş \pval-değerlerini yeniden üretir. Son hesap Welch sonucunu verir.
Ham veri bulunduğunda grup grafikleri, artıklar ve aykırı gözlemler ayrıca
incelenmelidir.
\end{softwarebox}

SPSS'te \emph{Analyze $>$ Compare Means $>$ One-Way ANOVA}; R'de klasik
analiz için \texttt{aov(y \textasciitilde{} group)}, Welch testi için
\texttt{oneway.test(y \textasciitilde{} group, var.equal=FALSE)} kullanılabilir.
Yazılım varsayılanlarının çoklu karşılaştırma ve eksik veri işlemleri
bakımından aynı olduğu varsayılmamalıdır.

\begin{outputguidebox}
Önce grup hacimleri, ortalamaları, standart sapmaları ve grafiklerini okuyun.
Ardından kullanılan genel testin klasik mi Welch mi olduğunu, pay ve payda
serbestlik derecelerini, $F$ ile \pval-değerini kontrol edin. Genel sonuçtan
sonra etki büyüklüğünü; gerekiyorsa düzeltilmiş ikili farkları ve eşzamanlı
güven aralıklarını raporlayın.
\end{outputguidebox}

\section{Raporlama ve araştırma tasarımı}

Klasik analiz için örnek rapor şöyledir:

``Başarı puanları öğretim yöntemlerine göre farklılaştı,
$F(2,27)=9{,}91$, $p<0{,}001$, $\eta^2=0{,}42$,
$\omega^2=0{,}37$. Tukey karşılaştırmaları C yönteminin A ve B'den daha
yüksek ortalamaya sahip olduğunu, A ile B farkının ise belirlenemediğini
gösterdi.''

Gruplara rastgele atama yapılmış ve uygulama dışında sistematik farklar
denetlenmişse yöntem etkisine ilişkin nedensel yorum güçlenir. Gözlemsel
gruplarda anlamlı ANOVA yalnız grup üyeliği ile yanıt arasında ilişki gösterir;
önceki başarı, seçim süreci veya başka karıştırıcılar farkı açıklayabilir.

\section{Çözümlü problemler}

\subsection*{Problem 1: Hipotez ve tasarım}

Dört farklı okul türündeki öğrencilerin aidiyet puanları karşılaştırılacaktır.
Yanıt değişkeni nicel ve öğrenciler yalnız bir okul türünde yer almaktadır.
Hipotezleri, uygun genel testi ve serbestlik derecelerini yazın. Toplam
örneklem hacmi 84'tür.

\textbf{Çözüm.} Dört bağımsız grup bulunduğundan tek yönlü bağımsız gruplar
ANOVA'sı uygundur. Hipotezler
\[
H_0:\mu_1=\mu_2=\mu_3=\mu_4,
\qquad H_1:\text{en az bir ortalama farklıdır}
\]
biçimindedir. Gruplar arası serbestlik derecesi $4-1=3$, hata serbestlik
derecesi $84-4=80$ ve toplam serbestlik derecesi $83$'tür.

\subsection*{Problem 2: Eksik ANOVA tablosunu tamamlama}

Üç grup ve toplam 36 gözlem için $SS_{\text{gruplar}}=240$,
$SS_{\text{hata}}=720$ verilmiştir. ANOVA tablosunu tamamlayın ve
$F_{0{,}95;2,33}=3{,}285$ kritik değeriyle karar verin.

\textbf{Çözüm.} Serbestlik dereceleri 2, 33 ve 35'tir.
\[
MS_{\text{gruplar}}=240/2=120,
\qquad MS_{\text{hata}}=720/33=21{,}818,
\]
\[
F=120/21{,}818=5{,}50.
\]
$5{,}50>3{,}285$ olduğundan yüzde 5 düzeyinde eşit ortalamalar hipotezi
reddedilir. Toplam kareler toplamı $240+720=960$'tır.

\subsection*{Problem 3: Etki büyüklükleri}

Problem 2'deki değerlerden $\eta^2$ ve $\omega^2$ hesaplayın.

\textbf{Çözüm.}
\[
\eta^2=240/960=0{,}25.
\]
$MS_{\text{hata}}=21{,}818$ olduğundan
\[
\omega^2
=\frac{240-2(21{,}818)}{960+21{,}818}
\approx0{,}200.
\]
Gruplama değişkeni örneklemde toplam değişkenliğin yüzde 25'iyle ilişkilidir;
omega-kare düzeltmesi bu oranı yaklaşık yüzde 20 olarak tahmin eder.

\subsection*{Problem 4: Genel test ve ikili karşılaştırma}

Bir ANOVA sonucunda $F(3,76)=4{,}20$, $p=0{,}008$ bulunmuştur. Altı ikili
karşılaştırmadan yalnız 1--4 çifti için Tukey düzeltilmiş $p=0{,}004$ ve
yüzde 95 eşzamanlı GA $[2{,}1;9{,}7]$ verilmiştir; diğer aralıklar sıfırı
içermektedir. Sonucu yorumlayın.

\textbf{Çözüm.} Genel test en az bir grup ortalamasının farklı olduğuna
ilişkin kanıt verir. Tukey sonuçları belirlenebilen farkın 1 ve 4. gruplar
arasında olduğunu; 1. grup ortalamasının 4. gruptan yaklaşık 2,1 ile 9,7 birim
daha yüksek olabileceğini gösterir. Diğer çiftlerde farkın belirlenememesi
ortalamaların kesin eşitliği anlamına gelmez.

\subsection*{Problem 5: Klasik ANOVA mı Welch ANOVA mı?}

Üç grubun hacimleri 8, 20 ve 40; standart sapmaları sırasıyla 15, 7 ve 4'tür.
Küçük grubun varyansı en büyüktür. Uygun genel analizi ve kontrol edilmesi
gereken noktaları açıklayın.

\textbf{Çözüm.} Grup hacimleri ve varyanslar belirgin biçimde dengesizdir;
klasik eşit varyanslı ANOVA'nın hata oranı etkilenebilir. Welch ANOVA uygun
birincil veya duyarlılık analizi olur. Grup kutu/nokta grafikleri, aykırı
değerler, bağımsızlık ve ölçüm koşulları incelenmelidir. İkili analizlerde de
Games--Howell gibi eşit varyans gerektirmeyen yöntem kullanılmalıdır.

\subsection*{Problem 6: Yazılım çıktısını bütüncül yorumlama}

Bir çıktıda grup ortalamaları 64, 66 ve 72; klasik ANOVA sonucu
$F(2,87)=6{,}40$, $p=0{,}003$, $\omega^2=0{,}11$ verilmiştir. Welch testi
$F_W(2,54{,}7)=5{,}92$, $p=0{,}005$ sonucunu vermekte; C grubunun kutu
grafiğinde iki aykırı değer adayı görülmektedir. Sonucu raporlayın.

\textbf{Çözüm.} Hem klasik hem Welch analizi ortalamaların eşit olmadığına
ilişkin benzer kanıt sunmaktadır. Örneklem ortalamaları en yüksek değerin C
grubunda olduğunu gösterir; fakat hangi çiftlerin farklı olduğu uygun çoklu
karşılaştırmayla belirlenmelidir. $\omega^2=0{,}11$ grup üyeliğiyle ilişkili
değişkenlik büyüklüğünü özetler. C grubundaki adaylar veri doğruluğu ve etki
bakımından incelenmeli; otomatik silinmemelidir. Sonuç, tasarım desteklemiyorsa
nedensel etki olarak sunulmamalıdır.

\section{R ile çözümlü uygulama}
\label{sec:r-anova}

\textbf{Soru.} Bölümdeki üç grubun $n=(10,10,10)$,
$\bar x=(70,75,82)$ ve $s=(6,7,5)$ özetlerinden klasik ANOVA'yı,
etki büyüklüklerini ve Tukey eşzamanlı fark aralıklarını hesaplayın.
Özetleri ham veri gibi \texttt{aov} komutuna vermeyin.

\begin{lstlisting}[style=prismR]
hacim <- c(10, 10, 10)
ortalama <- c(70, 75, 82)
standart_sapma <- c(6, 7, 5)
grup_sayisi <- length(hacim)
genel <- weighted.mean(ortalama, hacim)
ss_grup <- sum(hacim * (ortalama - genel)^2)
ss_hata <- sum((hacim - 1) * standart_sapma^2)
sd_grup <- grup_sayisi - 1
sd_hata <- sum(hacim) - grup_sayisi
mse <- ss_hata / sd_hata
f_degeri <- (ss_grup / sd_grup) / mse
c(F = f_degeri,
  p = pf(f_degeri, sd_grup, sd_hata, lower.tail = FALSE),
  eta2 = ss_grup / (ss_grup + ss_hata),
  omega2 = (ss_grup - sd_grup * mse) /
    (ss_grup + ss_hata + mse))
ciftler <- combn(seq_along(hacim), 2)
kritik <- qtukey(0.95, nmeans = grup_sayisi, df = sd_hata)
tukey <- t(apply(ciftler, 2, function(cift) {
  fark <- ortalama[cift[1]] - ortalama[cift[2]]
  se <- sqrt(mse / 2 * sum(1 / hacim[cift]))
  c(fark = fark, alt = fark - kritik * se,
    ust = fark + kritik * se,
    p_duzeltilmis = ptukey(abs(fark) / se,
      nmeans = grup_sayisi, df = sd_hata, lower.tail = FALSE))
}))
rownames(tukey) <- c("A-B", "A-C", "B-C")
tukey
\end{lstlisting}

\begin{outputguidebox}
\textbf{Çözüm ve kontrol değerleri.} Genel ortalama 75,6667;
$SS_{\text{grup}}=726{,}6667$, $SS_{\text{hata}}=990$ ve
$MSE=36{,}6667$'dir. $F(2,27)=9{,}9091$, $p=0{,}0005927$,
$\eta^2=0{,}4233$, $\omega^2=0{,}3726$ bulunur. Tukey aralıklarının
yarı genişliği 6,7143'tür. A--B aralığı $[-11{,}7143;1{,}7143]$,
A--C aralığı $[-18{,}7143;-5{,}2857]$, B--C aralığı
$[-13{,}7143;-0{,}2857]$'dir. Son iki aralık sıfırı dışlar.
\end{outputguidebox}

\textbf{Yorum.} Klasik ANOVA ve bu Tukey karşılaştırmaları bağımsız
gözlemler, grup içi yaklaşık normallik ve ortak hata varyansı modeline dayanır.
Genel testin anlamlılığı, üç çiftin hepsinin farklı olduğunu söylemez.
Ham veri bulunduğunda \texttt{aov(puan \textasciitilde{} grup, data = veri)}
ve ardından \texttt{TukeyHSD} kullanılabilir; grup değişkeni faktör olmalıdır.

\subsection*{Eşit olmayan varyanslarda ek çözüm}
Bölümdeki Welch örneğinin özetleriyle aynı düzeltmeyi doğrudan hesaplayın.

\begin{lstlisting}[style=prismR]
hacim <- c(12, 20, 8)
ortalama <- c(70, 76, 84)
varyans <- c(3, 10, 5)^2
grup_sayisi <- length(hacim)
agirlik <- hacim / varyans
pay <- agirlik / sum(agirlik)
merkez <- sum(pay * ortalama)
duzeltme <- sum((1 - pay)^2 / (hacim - 1))
f_welch <- (sum(agirlik * (ortalama - merkez)^2) /
  (grup_sayisi - 1)) /
  (1 + 2 * (grup_sayisi - 2) * duzeltme /
     (grup_sayisi^2 - 1))
sd_payda <- (grup_sayisi^2 - 1) / (3 * duzeltme)
c(F = f_welch, sd_pay = grup_sayisi - 1,
  sd_payda = sd_payda,
  p = pf(f_welch, grup_sayisi - 1, sd_payda,
         lower.tail = FALSE))
\end{lstlisting}

\textbf{Çözüm.} Yaklaşık $F=25{,}3265$, pay serbestlik derecesi 2 ve
payda serbestlik derecesi 18,1621 bulunur; $p=0{,}000005578$'dir. Ham veride
\texttt{oneway.test(puan \textasciitilde{} grup, data = veri, var.equal = FALSE)}
aynı yöntemi uygular. Klasik ANOVA'nın ortak MSE'siyle hesaplanan Tukey
aralıklarını bu eşit olmayan varyans örneğine taşımayın.

\textbf{Pekiştirme ve yanıt.} Klasik örnekte A--B farkı $-5$ olsa da
aralık sıfırı içerir. Bu, iki yöntemin eşdeğer olduğunun kanıtı değildir;
yüzde 5 aile düzeyinde fark gösterilemediğini söyler.

\section{Bölüm özeti}

\begin{itemize}
  \item Tek yönlü ANOVA, bir nicel yanıtın üç veya daha fazla bağımsız grup ortalamasını tek genel testle karşılaştırır.
  \item $F$ istatistiği gruplar arası ortalama kareyi grup içi ortalama kareye oranlar.
  \item Anlamlı genel test en az bir ortalamanın farklı olduğunu gösterir; farkın yerini tek başına göstermez.
  \item Eta-kare ve omega-kare genel farkın büyüklüğünü, çoklu karşılaştırmalar ise hangi grup çiftlerinin ayrıldığını açıklar.
  \item Varyanslar ve grup hacimleri belirgin biçimde farklıysa Welch ANOVA düşünülmelidir.
  \item Bağımsızlık, örnekleme ve atama süreci istatistiksel testten önce gelir; ANOVA tek başına nedensellik sağlamaz.
\end{itemize}

\section*{Alıştırmalar}

\begin{enumerate}
  \item Tek yönlü ANOVA'nın yanıtladığı araştırma sorusuna eğitim alanından bir örnek verin.
  \item Üç grup için $H_0$ ve $H_1$ hipotezlerini yazın.
  \item Neden bütün grup çiftlerine düzeltmesiz \tstat testleri uygulanmaması gerektiğini açıklayın.
  \item $k=5$ ve $N=120$ için gruplar arası, hata ve toplam serbestlik derecelerini bulun.
  \item $SS_{\text{gruplar}}=180$, $SS_{\text{hata}}=540$, $k=4$ ve $N=40$ için ortalama kareleri ve $F$ değerini hesaplayın.
  \item Önceki soruda $F_{0{,}95;3,36}=2{,}866$ ise test kararını yazın.
  \item $SS_{\text{gruplar}}=180$ ve $SS_{\text{toplam}}=720$ için $\eta^2$ değerini hesaplayın.
  \item Eta-kare ile omega-karenin neden aynı olmayabileceğini açıklayın.
  \item Anlamlı genel $F$ testinden sonra Tukey HSD'nin rolünü belirtin.
  \item Tukey fark aralığı $[-1{,}2;4{,}8]$ olan bir grup çiftini yorumlayın.
  \item Grup hacimleri 10, 10 ve 10; standart sapmaları 5, 6 ve 5 olduğunda klasik ANOVA'nın varyans koşulunu nitel olarak değerlendirin.
  \item Grup hacimleri 8, 20 ve 50; standart sapmaları 14, 7 ve 3 olduğunda Welch ANOVA'nın neden düşünülmesi gerektiğini açıklayın.
  \item Aynı öğrencilerin üç zamanda ölçüldüğü bir çalışmada bu bölümdeki bağımsız gruplar ANOVA'sının neden uygun olmadığını yazın.
  \item Kruskal--Wallis sonucunun hangi koşullarda doğrudan medyan farkı diye yorumlanamayacağını açıklayın.
  \item $F(2,57)=1{,}20$, $p=0{,}309$ sonucunu güvenli bir dille yorumlayın.
  \item İstatistiksel olarak anlamlı ANOVA sonucunun neden tek başına nedensel yöntem etkisi göstermediğini açıklayın.
\end{enumerate}